\section{Introduction}
\label{sec:intro}

Counting how many times an action repeats---push-ups in a workout, strokes of a
saw, turns of a wheel---is a basic form of temporal understanding with uses in
fitness, rehabilitation and industrial monitoring. It needs two abilities that
are rarely found in one model. The first is \emph{temporal precision at every
tempo}: on RepCount-A~\cite{hu2022transrac} the median cycle of a video ranges
from 0.85\,s to 4.75\,s (5th--95th percentile), and in one video out of five the
cycle length varies by more than 25\% as people accelerate, tire or pause. The
second is \emph{knowing what one repetition is}: a push-up is a down- and an
up-stroke, a cello bow stroke is one back-and-forth. The visual signal of many
actions repeats every half cycle, and which period counts is a convention of the
action, not a property of the motion.

Dedicated counters have the first ability and lack the second. They regress a
density or a period from video features at one temporal sampling rate, and their
largest errors are counts that are off by a factor of about two on symmetric
motions (\cref{sec:analysis}). Video multimodal large language models (MLLMs)
have the second ability and lack the first: asked how often an action repeats, a
recent 8B MLLM names the action and describes one repetition correctly, but its
counts are far off (zero-shot MAE 0.43--0.55 on four benchmarks). CountLLM~\cite{yao2025countllm}
fine-tunes an LLM to output the count as text, which ties both abilities to the
language model.

We present \method, a counter that takes temporal precision from a frozen
video encoder and semantics from a frozen MLLM. A \emph{multi-rate token
pyramid} encodes the video at three frame strides and two depths of the encoder
on a shared frame grid, and a decoder turns it into an event density whose sum is
the count. A frozen video MLLM reads the video once, is asked what is repeated
and what one repetition consists of, and its hidden state modulates the decoder;
the MLLM is never asked for the number. Two properties follow from this design.
First, because all pyramid levels live on one grid with stride-proportional clip
hops, \emph{shifting the pyramid by one level gives the exact encoder output of
the video played twice as fast}: tempo augmentation and a second test-time view
come for free, without re-encoding. Second, semantic conditioning lets one
counter learn from several datasets at once. Training on the union of
RepCount-A, Countix and UCFRep fails for an unconditioned counter---it improves
UCFRep but degrades RepCount-A, whose counting conventions are overridden by the
much larger Countix---and works once the decoder is told what is being repeated,
together with a learned dataset token and balanced sampling.

\TODO{results summary: RepCount-A, UCFRep, Countix test; QUVA zero-shot.}

Our contributions are:
\begin{itemize}
  \setlength\itemsep{0pt}
  \setlength\parskip{0pt}
  \item A division of labour for repetition counting in which a frozen video MLLM
  supplies what to count and a density decoder on a multi-rate token pyramid
  supplies how many, with an analysis of why each alone falls short.
  \item The pyramid-shift property, which turns the token pyramid into exact
  tempo-changed views for training and inference at no encoding cost.
  \item Semantically conditioned joint training across RAC datasets, with
  component ablations over five seeds and state-of-the-art results \TODO{claim
  after test runs}; code, splits, features and per-video predictions will be
  released.
\end{itemize}
